% CVaR / Expected Shortfall — research-note PDF.
%
% PROSE AND CODE ARE VERBATIM from the website article:
%   site/src/content/articles/cvar-expected-shortfall.tsx
% Like the VaR note, this article interpolates its computed results directly
% into the prose. Every interpolated value here is the one the site renders,
% read from data/cvar-expected-shortfall.ts. Code cells are inline in the
% article rather than in a -code.ts module.
%
% CHARTS generated by figures/make_cvar_expected_shortfall.py from that module.
\documentclass{pyportfolios-note}

\noteshorttitle{CVaR / Expected Shortfall}

\begin{document}
\thispagestyle{notecover}

\notemasthead{Quantitative Insights}{Tutorial}{Q3 2026}{}

\notetitle{CVaR / Expected\\Shortfall}

\notedeck{VaR tells you where the tail begins. Expected shortfall tells you what
lives inside it — which is where portfolios actually die.}

\notelede{Value-at-Risk tells you where the tail begins. It says nothing about
what lives inside it — and inside the tail is where portfolios actually die. Two
books can report identical 99\% VaR yet differ threefold in what they lose once
the threshold breaks; for fat-tailed assets like high-yield credit, that
difference is the whole risk. Expected shortfall (CVaR) fixes the blind spot by
averaging the tail instead of pointing at its door, which is why Basel made it
the regulatory standard. We estimate both, historically and with a Student-t, on
4,462 days of HYG and VWO — a sample that deliberately includes the crisis these
measures were built for.}

\notecard
  {Risk Management}
  {NumPy · SciPy · Pandas · Riskfolio-Lib}
  {HYG + VWO}
  {Apr 2007 – Dec 2024}
  {Setting tail-risk limits that survive a crisis — and pass a regulator}

\begin{multicols}{2}
\begin{notepipeline}
\item Load HYG + VWO daily closes, 2007–2024 — the GFC included
\item Build two synthetic books with identical 99\% VaR, very different tails
\item Estimate 97.5\% / 99\% VaR \& CVaR — historical and Student-t (Acerbi's
      formula)
\item Test subadditivity: CVaR always passes, VaR fails in crisis-year samples
\item Zoom into HYG's GFC tail against its full-sample 99\% VaR
\item Spend the same 4\% tail budget through rm=``MV'' vs rm=``CVaR''
      (Riskfolio-Lib)
\end{notepipeline}
\end{multicols}

% =========================================== 1 What VaR cannot see ===========
\noteneedroom{190bp}
\notesection{1.\notenumsep What VaR cannot see}

\begin{multicols}{2}
VaR at confidence α is a \textbf{quantile}: the smallest loss exceeded on only
the worst (1−α) of days. Read that definition again and notice what it does not
say — nothing about \textbf{how much} you lose on those days. VaR is a
threshold, not an average, so it is structurally blind to everything beyond
itself. A ten-line thought experiment makes the blindness concrete: build two
1,000-day P\&L histories that agree on 990 benign days \textbf{and} on the day
that sets the 99\% quantile, but whose ten worst days differ by a factor of six.

Both books report a 99\% VaR of 2.40\%, so a risk report built on VaR alone
calls them equally risky. Yet book A's CVaR is 2.60\% against book B's 7.26\%,
with worst days of −2.60\% and −15.00\% — identical VaR, 2.8× the expected tail
loss. CVaR — the \textbf{mean} loss conditional on breaching the VaR quantile —
separates the two books instantly, because it averages over the whole tail
instead of reading one point at its edge.
\end{multicols}

\notecallout{Why practitioners care}{Basel's Fundamental Review of the Trading
Book (FRTB) replaced 99\% VaR with 97.5\% expected shortfall as the capital
measure for trading desks — precisely because VaR ignores tail severity and can
punish diversification. The 97.5\% level was chosen so that, for a normal
distribution, ES roughly matches the old 99\% VaR — but unlike VaR it keeps
paying attention when the tail turns out fatter than normal.}

% ============================================== 2 Two fat-tailed assets ======
\noteneedroom{190bp}
\notesection{2.\notenumsep Two fat-tailed assets}

Now the real thing. HYG (iShares high-yield corporate bond ETF) and VWO
(Vanguard emerging-markets equity) are chosen because they are risky in opposite
ways. EM equity is honestly volatile — it looks risky and it is. Credit is the
trap: it collects small, steady coupons in calm markets and hands them back all
at once in a crisis. On paper HYG's 11.2\% annualised vol looks far safer than
VWO's 27.1\% — until you notice that the quiet series hides a worst day of
−8.10\%. The aligned sample runs 2007-04-11 (HYG's listing) to 2024-12-31. A
Student-t fitted to HYG's daily returns lands at ν ≈ 1.98 degrees of freedom —
below 2, which means the fitted distribution does not even possess a finite
variance. Read that as a diagnostic as much as an estimate: an unconditional iid
fit has nowhere to put 2008's volatility clustering except the tail parameter, so
it buys realism at the extremes by overstating how wild a \textbf{typical} day
is. (Condition on a GARCH filter and the residual df comes out higher; the
unconditional fit is the honest worst case.)

\notefigure{figures/cvar-expected-shortfall/fig1-hyg-hist.png}
  {HYG daily returns 2007 → 2024 with fitted Student-t (ν = 1.98) — VaR marks the door, CVaR the room behind it}
  {The horizontal gap between the two thresholds is everything VaR does not
   price}
  {Source: Author calculations, from the article's computed results.}

The horizontal gap between the two dashed lines — 2.09\% to 3.37\% — is
everything VaR does not price, and HYG's gap is the widest of the three books we
measure.

% ============================= 3 Historical and Student-t estimates ==========
\noteneedroom{190bp}
\notesection{3.\notenumsep Historical and Student-t estimates}

Estimating CVaR comes down to two philosophies. The historical estimator trusts
the data: sort, take the quantile, average everything beyond it. The parametric
route trusts the model: fit a Student-t and use \textbf{Acerbi's closed-form
expected shortfall} — the analytic mean of the t's tail — which can extrapolate
to days worse than anything the sample has actually seen:

\begin{notecode}
def var_hist(x, a):
    return -np.quantile(x, 1 - a)

def cvar_hist(x, a):
    q = np.quantile(x, 1 - a)
    return -x[x <= q].mean()

def cvar_t(params, a):                    # Acerbi & Tasche closed form
    nu, loc, scale = params
    p = 1 - a
    xp = stats.t.ppf(p, nu)
    return -loc + scale * stats.t.pdf(xp, nu) * (nu + xp**2) / ((nu - 1) * p)
\end{notecode}

\begin{notetable}{@{}p{250bp}>{\raggedleft\arraybackslash}p{115bp}>{\raggedleft\arraybackslash}p{115bp}@{}}
\thcell{HYG estimate} & \thcell{97.5\%} & \thcell{99\%} \\
\theadrule
Historical VaR & 1.28\% & 2.09\% \\ \rowrule
Historical CVaR & 2.29\% & 3.37\% \\ \rowrule
Student-t VaR (ν = 1.98) & 1.21\% & 1.99\% \\ \rowrule
Student-t CVaR (Acerbi) & 2.54\% & 4.09\% \\
\end{notetable}

Three things to read off that table:

\begin{multicols}{2}
\begin{notebullets}
\item \textbf{The FRTB calibration at work} — HYG's 97.5\% CVaR (2.29\%) sits
      close to its 99\% VaR (2.09\%): similar magnitude, but the CVaR number
      keeps growing when the tail does.
\item \textbf{The t extrapolates past the sample} — its CVaR (4.09\%) exceeds
      the historical one because with ν ≈ 1.98 the fitted tail expects days worse
      than any yet observed.
\item \textbf{The small-sample caveat} — at 99\% the historical CVaR is the mean
      of just 45 observations. The estimator with the best theoretical properties
      stands on the fewest data points, which is the practical argument for
      fitting a parametric tail and letting it extrapolate.
\end{notebullets}
\end{multicols}

\notefigure{figures/cvar-expected-shortfall/fig2-var-cvar.png}
  {99\% VaR vs 99\% CVaR, historical — per asset and for the daily-rebalanced 50/50}
  {For every book the CVaR bar runs roughly one and a half times the VaR bar}
  {Source: Author calculations, from the article's computed results.}

\begin{notetable}{@{}p{72bp}>{\raggedleft\arraybackslash}p{72bp}>{\raggedleft\arraybackslash}p{72bp}>{\raggedleft\arraybackslash}p{78bp}>{\raggedleft\arraybackslash}p{72bp}p{102bp}@{}}
\thcell{Asset} & \thcell{Ann. vol} & \thcell{99\% VaR} & \thcell{99\% CVaR} & \thcell{CVaR/VaR} & \thcell{Worst day} \\
\theadrule
HYG & 11.2\% & 2.09\% & 3.37\% & 1.62× & −8.10\% (2008-09-29) \\ \rowrule
VWO & 27.1\% & 4.86\% & 7.15\% & 1.47× & −15.34\% (2008-10-15) \\ \rowrule
50/50 & 17.5\% & 3.26\% & 4.92\% & 1.51× & −9.85\% (2008-10-15) \\
\end{notetable}

Read the first and fifth columns together. HYG's volatility is well under half of
VWO's, yet its CVaR/VaR ratio of 1.62× is the \textbf{highest} in the table —
credit's deceptively quiet variance hides the most disproportionate tail. Any
tool that ranks these assets by σ alone has the risk ordering half wrong.

% ================================= 4 Subadditivity — the coherence test ======
\noteneedroom{190bp}
\notesection{4.\notenumsep Subadditivity — the coherence test}

\begin{multicols}{2}
In 1999, Artzner, Delbaen, Eber and Heath wrote down four axioms any sane risk
measure should satisfy. VaR fails one — and it is the one your intuition cares
most about: \textbf{subadditivity}, ρ(A+B) ≤ ρ(A) + ρ(B), the rule that
diversification must never create risk. The classic counterexample needs only two
independent bonds, each defaulting with probability 0.7\%. Held alone, each has
\textbf{zero} 99\% VaR: a 0.7\% loss probability hides entirely below the 1\%
threshold, so VaR literally cannot see it. Mix them 50/50 and the book loses on
1.4\% of scenarios — above 1\% — so the diversified book has strictly
\textbf{positive} 99\% VaR. Diversifying ``created'' risk, says VaR. The risk was
there all along, of course. VaR was hiding it.

It is not just a parlour trick. On the full HYG/VWO sample the 50/50 portfolio
behaves — 99\% VaR of 3.26\% against a weighted blend of 3.47\%. But scan
calendar-year subsamples across confidence levels and historical VaR breaks
subadditivity 64 times; the worst case is 2009 at α = 98.8\%, where the
portfolio's VaR of 4.48\% exceeds the blend's 3.95\%. CVaR, by construction —
Rockafellar \& Uryasev's convex formulation makes this explicit — never violates:
4.92\% for the portfolio versus 5.26\% for the blend at 99\%, and it passes at
every year and level where VaR fails.
\end{multicols}

\notecallout{Note}{Subadditivity is what makes a risk measure \textbf{aggregable}:
desk CVaRs can be summed into a conservative firm-wide bound, and CVaR limits can
be pushed through a convex optimiser. Neither is true of VaR — a VaR-constrained
optimisation is non-convex and can reward concentration.}

% ========================================= 5 The GFC, seen from the tail =====
\noteneedroom{190bp}
\notesection{5.\notenumsep The GFC, seen from the tail}

HYG's full-sample 99\% VaR is 2.09\%. Now watch the crisis ignore it: from
September 2008 through March 2009, HYG breached that threshold on 24 of 146
trading days — a 1\% tail arriving at 24× its expected frequency of
\textasciitilde1.5 days. The worst prints came fast: −8.10\% on 2008-09-29,
−6.67\% on 2008-10-10, −6.01\% on 2008-09-17 — every one of them multiples of the
VaR line, in an ETF marketed as a bond fund.

\notefigure{figures/cvar-expected-shortfall/fig3-gfc.png}
  {HYG cumulative return, 2007-04-11 → 2010 — the GFC drawdown that a 2.09\% VaR never described}
  {Roughly flat, then a collapse to about a third below its start into early
   2009, then recovery}
  {Source: Author calculations, from the article's computed results.}

At the trough a dollar invested at HYG's listing was worth \$0.69. VaR answered
``how often?'' — and even that answer failed under regime change. CVaR at least
asks the question that mattered in 2008: \textbf{how bad is it when it happens?}

% =============================== 6 Spending a tail budget: CVaR vs MV ========
\noteneedroom{190bp}
\notesection{6.\notenumsep Spending a tail budget: CVaR vs MV}

So far CVaR has been the better thermometer. The sharper question is whether it
changes what you \textbf{do}. With two assets, pure risk-minimising is degenerate
here — volatility and tail agree that HYG is the quieter asset, so both
\texttt{rm="MV"} and \texttt{rm="CVaR"} corner at 100\% HYG. The measures diverge
the moment you \textbf{spend a risk budget}. Hand two desks the same mandate — an
expected loss on the worst 1\% of days of at most 4\% — and let each maximise
return against it. The MV desk speaks only volatility, so it must translate the
budget through a normal distribution (ES₉₉ = 2.665σ for a Gaussian, so σ ≤ 1.50\%
daily); the CVaR desk constrains the realised tail directly:

\begin{notecode}
import riskfolio as rp

p_mv = rp.Portfolio(returns=rets[["HYG", "VWO"]])
p_mv.assets_stats(method_mu="hist", method_cov="hist")
p_mv.upperdev = 0.015                    # 4% ES budget, normal-translated
w_mv = p_mv.optimization(model="Classic", rm="MV", obj="MaxRet", hist=True)

p_cv = rp.Portfolio(returns=rets[["HYG", "VWO"]])
p_cv.assets_stats(method_mu="hist", method_cov="hist")
p_cv.alpha = 0.01
p_cv.upperCVaR = 0.04                      # the tail budget itself
w_cv = p_cv.optimization(model="Classic", rm="CVaR", obj="MaxRet", hist=True)
\end{notecode}

\notefigure{figures/cvar-expected-shortfall/fig4-budget.png}
  {Same 4\% tail budget, two allocations — how each desk spends it}
  {The MV desk puts 83.8\% in VWO where the CVaR desk allows only 24.7\%}
  {Source: Author calculations, from the article's computed results.}

\begin{notetable}{@{}p{198bp}>{\raggedleft\arraybackslash}p{90bp}>{\raggedleft\arraybackslash}p{80bp}p{128.8bp}@{}}
\thcell{Desk} & \thcell{Ann. return} & \thcell{Ann. vol} & \thcell{Realised 99\% CVaR} \\
\theadrule
MV (vol-translated budget) & 6.6\% & 23.8\% & 6.37\% — budget was 4\% \\ \rowrule
CVaR (direct) & 5.7\% & 13.6\% & 3.99\% \\
\end{notetable}

Same stated risk appetite, radically different books: 83.8\% VWO through the
variance lens versus 24.7\% through the CVaR lens. Because both assets' tails are
fatter than the Gaussian used in the translation, the MV desk's realised 99\%
CVaR of 6.37\% overshoots its own 4\% mandate by roughly 59\% — the tail it was
told to cap is exactly what its risk measure could not see.

\notecallout{Practitioner take}{Use CVaR where the loss distribution is
asymmetric or fat-tailed — credit, options, carry, anything with a ``smile now,
cry later'' profile. HYG is the canonical trap: its vol whispers safety while its
tail shouts the opposite. Practically: quote 97.5\% ES next to any 99\% VaR
(FRTB's own calibration), fit a Student-t rather than trusting the empirical tail
alone when ν comes out below \textasciitilde4, and put the CVaR constraint
\textbf{inside} the optimiser — Rockafellar–Uryasev makes it a linear program, so
there is no computational excuse. The one honest cost of the switch: ES is harder
to \textbf{backtest} than VaR. A quantile is directly falsifiable by counting
breaches; an expected shortfall is not elicitable on its own, only jointly with
its VaR — which is why desk validation (Acerbi–Székely) and FRTB's own
backtesting still run on VaR exceptions at two confidence levels, even though the
capital number is ES.}

\noteneedroom{200bp}
\begin{notereferences}
\item Artzner, P., Delbaen, F., Eber, J.-M. \& Heath, D. (1999). Coherent
      Measures of Risk. \textit{Mathematical Finance} 9(3), 203–228.
\item Acerbi, C. \& Tasche, D. (2002). On the coherence of expected shortfall.
      \textit{Journal of Banking \& Finance} 26(7), 1487–1503.
\item Rockafellar, R.T. \& Uryasev, S. (2000). Optimization of Conditional
      Value-at-Risk. \textit{Journal of Risk} 2(3), 21–41.
\item Acerbi, C. \& Székely, B. (2014). Back-testing Expected Shortfall.
      \textit{Risk Magazine}, December 2014.
\item Basel Committee on Banking Supervision (2019). Minimum capital
      requirements for market risk (FRTB). Bank for International Settlements.
\item Companion notebook: \texttt{cvar-expected-shortfall.ipynb} — reproduces
      every figure from raw data (fully deterministic; no simulation).
\end{notereferences}

\end{document}
